% ============================================================================
% Fichier  : annexes/Ann_E_bibliographie.tex
% Titre    : Bibliographie Commentée
% Version  : 1.0 -- novembre 2024
% ============================================================================

\chapter{Bibliographie Commentée}
\label{ann:bibliographie}

Les références complètes figurent dans la bibliographie en fin d'ouvrage.
Cette annexe indique, partie par partie, par où commencer.

\section*{Partie I. Fondements}

E.~Casey~\cite{casey2004} pose les notions de trace et de preuve numérique
avec de nombreux exemples judiciaires. Le guide NIST SP~800-86~\cite{nist80086},
court et gratuit, est la meilleure entrée dans les phases d'une enquête.
R.~Heuer~\cite{heuer1999} reste la référence sur les biais de l'analyste et
l'analyse des hypothèses concurrentes. L'ISO/IEC~27037~\cite{iso27037_2012}
et l'ISO/IEC~27043~\cite{iso27043_2015} donnent le cadre normatif.

\section*{Partie II. Méthode}

Pour le disque, B.~Carrier~\cite{carrier2005} explique les systèmes de
fichiers de l'intérieur~; on le lit avec Autopsy ouvert. Pour la mémoire,
M.~H.~Ligh et ses coauteurs~\cite{ligh2014memory} couvrent Windows, Linux et
macOS~; la documentation de Volatility~3~\cite{volatility3} met à jour les
commandes. Pour le réseau, R.~Bejtlich~\cite{bejtlich2013nsm} pense la
supervision, C.~Sanders~\cite{sanders2017packet} apprend à lire les paquets.
La RFC~3227~\cite{rfc3227} tient en quelques pages et se relit avant chaque
intervention.

\section*{Partie III. Cadre normatif et juridique}

La loi n°~2010/012~\cite{loi2010012} se lit dans son texte officiel,
chapitre de procédure d'abord. Les Conventions de
Budapest~\cite{budapest2001} et de Malabo~\cite{ua2014malabo} éclairent
l'entraide.

\section*{Partie IV. Vulnérabilités}

Les spécifications du CVSS~4.0~\cite{first2023cvss4}, la documentation de
l'EPSS~\cite{firstepss}, le catalogue KEV~\cite{cisakev} et le guide
SSVC~\cite{cisa2022ssvc} sont gratuits et courts. Le NIST
SP~800-40~\cite{nist80040} traite la gestion des correctifs comme une
maintenance planifiée~; le NIST SP~800-115~\cite{nist800115} distingue
évaluation et test d'intrusion.

\section*{Partie V. Défis avancés}

L'ISO/IEC~27042~\cite{iso27042_2015} pour l'interprétation et le rapport~;
l'ISO/IEC~27035~\cite{iso27035} et le NIST SP~800-61~\cite{nist80061r2} pour
les enseignements tirés d'un incident.

\section*{Partie VI. Recherche}

Les deux articles fondateurs du programme de l'auditologie sont déposés sur
HAL~: « Vers l'auditologie »~\cite{minka2026auditologie} et l'article sur
l'instrument trichotomique~\cite{minka2026nek}. Les articles suivants du
programme sont en préparation. C.~E.~Shannon~\cite{shannon1948} reste la
source de la mesure d'entropie.
